\label{sec:related_work}

\textbf{Ontology-based IR and food ontologies.}
Classical ontology-based retrieval enriches document and query representations with explicit semantic structure. Vallet et al.~\cite{vallet2005ontology} combine keyword search with ontology annotations, Castells et al.~\cite{castells2007adaptation} adapt the vector-space model to ontology-based ranking, and Fern\'{a}ndez et al.~\cite{fernandez2011semantic} extend this line to larger-scale settings, collectively showing that concept-level representations reduce vocabulary mismatch and improve retrieval precision. In the food domain, FoodOn~\cite{dooley2018foodon} provides a large harmonized food ontology for traceability and data integration, while FoodKG~\cite{haussmann2019foodkg} links recipes, ingredients, and nutrients in a knowledge graph used for explainable recommendation. Chen et al.~\cite{chen2021personalized} further formulate personalized food recommendation as constrained question answering over a food knowledge graph. However, these food-oriented studies target recommendation or question answering rather than ranking-based retrieval, and the classical ontology-IR studies predate dense retrieval. Neither line integrates ontology with modern neural retrieval pipelines, particularly for low-resource languages such as Vietnamese.

\textbf{Food retrieval under neural frameworks.}
Recent work increasingly treats food search as a retrieval problem
within neural pipelines. Early recipe retrieval studies focused on
cross-modal alignment between recipes and images
\cite{carvalho2018crossmodal, lien2020recipe}. More recent benchmarks
such as Recipe-MPR \cite{zhang2023recipempr} highlight the difficulty
of multi-attribute preference-based retrieval, while Hu et
al.~\cite{hu2024bridging} extend recipe retrieval to cross-cultural
settings through reformulation and reranking. KERL
\cite{mohbat2025kerl} further shows that structured food knowledge
remains useful even in LLM-augmented pipelines. In the broader IR
community, dense retrieval has become the dominant paradigm for
semantic matching \cite{karpukhin2020dense, thakur2021beir}, and
recent graph-aware retrieval frameworks demonstrate that explicit
relational structure can complement purely embedding-based retrieval
on compositional queries \cite{he2024g, hipporag2024, kg2rag2025}.

\textbf{Ingredient substitution and dish similarity.}
Ingredient substitution has been studied through both knowledge-graph and learning-based approaches. Shirai et al.~\cite{shirai2021identifying} identify substitutions using a food knowledge graph, while Fatemi et al.~\cite{fatemi2023learning} learn substitution patterns from recipe corpora. For dish similarity, structured food relations have also been modeled through food-chemical or graph-based representations such as FlavorGraph \cite{park2021flavorgraph}. However, lightweight ingredient-set baselines often compute dish similarity from flat ingredient overlap, without distinguishing substitutions within the same semantic class from cross-class differences.

\textbf{Position of the present work.}
Our work bridges these three lines by: (1)~constructing a Vietnamese-specific ontology from a 10K-dish dataset, (2)~integrating it into a dense-retrieval pipeline as a pre-embedding enrichment mechanism for three food retrieval tasks, and (3)~providing an explicit ablation that isolates the contribution of hierarchy, typed relations, and inference rules.
